\section{转动惯量与惯性积的计算公式}

下面给出矩形块 $B$ 关于质心 $B^{*}$ 的转动惯量与惯性积的一般计算公式。

\subsubsection*{（一）转动惯量与惯性积的积分定义}
对刚体上任一质点 $P\in B$，设其相对质心 $B^{*}$ 的位矢为 $\boldsymbol{r}=x_{1}\boldsymbol{b}_{1}+x_{2}\boldsymbol{b}_{2}+x_{3}\boldsymbol{b}_{3}$，则惯性并矢 $\mathbf{I}$ 在 $B$ 系的分量为
\begin{equation}
    \mathbf{I}_{ij|B}\;\triangleq\;\int_{B}\rho\,x_{i}x_{j}\,dV \qquad (i,j=1,2,3).
\end{equation}
对角元 $\mathbf{I}_{ii|B}$ 称为\textbf{标量转动惯量}，非对角元 $\mathbf{I}_{ij|B}\,(i\neq j)$ 称为\textbf{惯性积（积矩）}。

\subsubsection{（二）标量转动惯量的计算}
设矩形块 $B$ 的边长分别为 $2a,2b,2c$（沿 $\boldsymbol{b}_{1},\boldsymbol{b}_{2},\boldsymbol{b}_{3}$ 方向），密度 $\rho$ 均匀。则绕 $\boldsymbol{b}_{1}$ 轴的转动惯量是质点到 $\boldsymbol{b}_{1}$ 轴的\textbf{垂直距离的平方}乘以密度的体积分：
\begin{equation}
\begin{aligned}
\mathbf{I}_{11|B}
  &= \int_{B}\rho\,(x_{2}^{2}+x_{3}^{2})\,dV \\
  &= \rho\int_{-c}^{c}\!\!\int_{-b}^{b}\!\!\int_{-a}^{a}(y^{2}+z^{2})\,dx\,dy\,dz \\
  &= \frac{\rho\cdot 2a\cdot 2b\cdot 2c}{12}\,(2b^{2}+2c^{2}) \\
  &= \frac{m}{12}(b^{2}+c^{2}).
\end{aligned}
\end{equation}
类似地（轮换 $a,b,c$）：
\begin{equation}
    \mathbf{I}_{22|B}=\frac{m}{12}(a^{2}+c^{2}),\qquad
    \mathbf{I}_{33|B}=\frac{m}{12}(a^{2}+b^{2}).
\end{equation}

$\frac{1}{12}$ 这一系数来源于：在 $[-h/2,h/2]$ 上 $\displaystyle\frac{1}{h}\int_{-h/2}^{h/2}y^{2}\,dy=\frac{h^{2}}{12}$，即均匀分布的均方距离。

\subsubsection{（三）惯性积的计算}
惯性积的定义为
\begin{equation}
    \mathbf{I}_{ij|B}=\int_{B}\rho\,x_{i}x_{j}\,dV \qquad (i\neq j).
\end{equation}
对矩形块：
\begin{equation}
\begin{aligned}
\mathbf{I}_{12|B}
  &= \int_{B}\rho\,x_{1}x_{2}\,dV \\
  &= \rho\int_{-c}^{c}\!\!\int_{-b}^{b}\!\!\int_{-a}^{a}xy\,dx\,dy\,dz.
\end{aligned}
\end{equation}
注意到被积函数 $xy$ 关于 $x=0$ 是奇函数，在 $[-a,a]$ 上的积分为零：
\begin{equation}
    \int_{-a}^{a}x\,dx = 0 \quad\Longrightarrow\quad \mathbf{I}_{12|B}=0.
\end{equation}
由对称性同理可得
\begin{equation}
    \mathbf{I}_{13|B}=\mathbf{I}_{23|B}=0.
\end{equation}
这就是 $\boldsymbol{b}_{1},\boldsymbol{b}_{2},\boldsymbol{b}_{3}$ 之所以被称为惯性主轴的根本原因。

\subsubsection{（四）$\mathbf{I}_{ij|B}$ 的对称性}
由定义 $\mathbf{I}_{ij|B}=\int\rho\,x_{i}x_{j}\,dV$ 显然有
\begin{equation}
    \mathbf{I}_{ij|B}=\mathbf{I}_{ji|B},
\end{equation}
即惯性张量是对称张量：对角元（转动惯量）只需计算 3 个，非对角元（惯性积）只需计算 3 个，共 6 个独立分量。

\subsubsection{（五）代入矩形块的具体尺寸}
将反推得到的 $a=4L,\,b=12L,\,c=3L$ 代入上面的一般公式，立即得到上面给出的三个主对角元
\[
\mathbf{I}_{11|B}=\frac{m}{12}(12^{2}+3^{2})L^{2}=\frac{153}{12}mL^{2},\quad
\mathbf{I}_{22|B}=\frac{m}{12}(3^{2}+4^{2})L^{2}=\frac{25}{12}mL^{2},\quad
\mathbf{I}_{33|B}=\frac{m}{12}(4^{2}+12^{2})L^{2}=\frac{160}{12}mL^{2}.
\]
